\documentclass[10pt,twocolumn]{article}
\usepackage[margin=0.6in]{geometry}
\usepackage{amsmath, amssymb, amsthm, enumitem, booktabs, array}

\setlength{\parindent}{0pt}
\setlength{\parskip}{4pt}
\pagestyle{empty}

\newcommand{\sect}[1]{\medskip\noindent\textbf{\large #1}\medskip\hrule\smallskip}
\newcommand{\subsect}[1]{\smallskip\noindent\textbf{#1}\par}
\DeclareMathOperator{\rank}{rank}
\DeclareMathOperator{\nullity}{nullity}
\DeclareMathOperator{\tr}{tr}
\DeclareMathOperator{\Span}{span}
\DeclareMathOperator{\im}{im}
\DeclareMathOperator{\diag}{diag}

\begin{document}

\begin{center}
{\LARGE\bfseries GRE Math Subject Test}\\[2pt]
{\Large Linear Algebra Concept Sheet}
\end{center}

\vspace{-4pt}

%======================================================================
\sect{1. Vector Spaces}
%======================================================================

\subsect{Axioms}
A \textbf{vector space} $V$ over a field $F$ is an abelian group $(V,+)$ with scalar multiplication $F\times V\to V$ satisfying: associativity, distributivity (both kinds), and $1\cdot v=v$.

\subsect{Key Examples}
\begin{itemize}[nosep,leftmargin=*]
\item $\mathbb{R}^n$, $\mathbb{C}^n$ over their respective fields.
\item $M_{m\times n}(F)$: $m\times n$ matrices over $F$.
\item $F[x]$: polynomials over $F$ (infinite-dimensional).
\item $P_n(F)$: polynomials of degree $\le n$ ($\dim = n+1$).
\item $\{0\}$: the zero vector space ($\dim=0$).
\end{itemize}

\subsect{Subspaces}
$W\subseteq V$ is a subspace if:
\begin{itemize}[nosep,leftmargin=*]
\item $0\in W$, and
\item $W$ is closed under addition and scalar multiplication.
\end{itemize}
Intersection of subspaces is a subspace. Union generally is NOT.

\subsect{Span, Independence, Basis, Dimension}
\begin{itemize}[nosep,leftmargin=*]
\item $\Span(S)=$ all linear combinations of vectors in $S$.
\item $\{v_1,\ldots,v_k\}$ is \textbf{linearly independent} if $c_1v_1+\cdots+c_kv_k=0 \Rightarrow$ all $c_i=0$.
\item \textbf{Basis:} linearly independent spanning set.
\item $\dim V=$ number of vectors in any basis.
\item In $\mathbb{R}^n$: $\dim=n$.\ In $M_{m\times n}$: $\dim=mn$.\ In $P_n$: $\dim=n+1$.
\end{itemize}

\subsect{Dimension Formula}
$\dim(W_1+W_2)=\dim W_1+\dim W_2-\dim(W_1\cap W_2).$

%======================================================================
\sect{2. Linear Transformations}
%======================================================================

\subsect{Definition}
$T\colon V\to W$ is linear if $T(\alpha u+\beta v)=\alpha T(u)+\beta T(v)$.

\subsect{Kernel and Image}
\begin{itemize}[nosep,leftmargin=*]
\item $\ker T=\{v\in V:T(v)=0\}$ (subspace of $V$).
\item $\im T=\{T(v):v\in V\}$ (subspace of $W$).
\item $T$ injective $\iff$ $\ker T=\{0\}$.
\item $T$ surjective $\iff$ $\im T=W$.
\end{itemize}

\subsect{Rank--Nullity Theorem}
For $T\colon V\to W$ with $\dim V<\infty$:
\[
\dim V = \rank(T)+\nullity(T) = \dim(\im T)+\dim(\ker T).
\]

\subsect{Isomorphism}
$T$ is an isomorphism $\iff$ $T$ is bijective and linear.\\
$V\cong W$ as vector spaces $\iff$ $\dim V=\dim W$ (finite-dimensional case).

%======================================================================
\sect{3. Matrices \& Row Reduction}
%======================================================================

\subsect{Matrix of a Linear Map}
$T\colon V\to W$ with ordered bases $B_V$, $B_W$: the matrix $[T]_{B_V}^{B_W}$ has columns = coordinates of $T(v_j)$ in basis $B_W$.

\subsect{Row Echelon Form (REF)}
Apply elementary row operations to get REF/RREF.
\begin{itemize}[nosep,leftmargin=*]
\item $\rank(A)=$ number of pivots = number of nonzero rows in REF.
\item $\nullity(A)=n-\rank(A)$ (for $A$ with $n$ columns).
\item $Ax=b$ consistent $\iff$ $\rank(A)=\rank([A\mid b])$.
\item $Ax=0$ has nontrivial solutions $\iff$ $\rank(A)<n$.
\end{itemize}

\subsect{Invertibility (Equivalent Conditions)}
For $n\times n$ matrix $A$, the following are equivalent:
\begin{itemize}[nosep,leftmargin=*]
\item $A$ is invertible.
\item $\det A\neq 0$.
\item $\rank(A)=n$.
\item $\ker A=\{0\}$ (columns are linearly independent).
\item $Ax=b$ has a unique solution for every $b$.
\item $0$ is NOT an eigenvalue of $A$.
\item RREF of $A$ is $I_n$.
\end{itemize}

\subsect{Matrix Algebra}
\begin{itemize}[nosep,leftmargin=*]
\item $(AB)^T=B^TA^T$, \quad $(AB)^{-1}=B^{-1}A^{-1}$.
\item $\rank(AB)\le\min(\rank A,\rank B)$.
\item $\rank(A^T A)=\rank(A)$.
\end{itemize}

%======================================================================
\sect{4. Determinants}
%======================================================================

\subsect{Properties}
\begin{itemize}[nosep,leftmargin=*]
\item $\det(AB)=\det A\cdot\det B$.
\item $\det(A^T)=\det A$.
\item $\det(A^{-1})=1/\det A$.
\item $\det(cA)=c^n\det A$ for $n\times n$ matrix $A$.
\item Swapping two rows: $\det\to-\det$.
\item Row of zeros $\Rightarrow$ $\det=0$.
\item Two identical rows $\Rightarrow$ $\det=0$.
\item Adding a multiple of one row to another: $\det$ unchanged.
\end{itemize}

\subsect{Computation}
\begin{itemize}[nosep,leftmargin=*]
\item $2\times 2$: $\det\begin{pmatrix}a&b\\c&d\end{pmatrix}=ad-bc$.
\item $3\times 3$: cofactor expansion or Sarrus' rule.
\item Triangular/diagonal: $\det=$ product of diagonal entries.
\item Block triangular: $\det\begin{pmatrix}A&*\\0&D\end{pmatrix}=\det A\cdot\det D$.
\end{itemize}

\subsect{Geometric Meaning}
$|\det A|=$ volume scaling factor of the linear map.

%======================================================================
\sect{5. Eigenvalues \& Eigenvectors}
%======================================================================

\subsect{Definitions}
$Av=\lambda v$ ($v\neq 0$): $\lambda$ is an \textbf{eigenvalue}, $v$ is an \textbf{eigenvector}.\\
\textbf{Eigenspace:} $E_\lambda=\ker(A-\lambda I)$.

\subsect{Characteristic Polynomial}
$p(\lambda)=\det(A-\lambda I)$. Eigenvalues = roots of $p(\lambda)=0$.
\begin{itemize}[nosep,leftmargin=*]
\item Degree $n$ polynomial $\Rightarrow$ at most $n$ distinct eigenvalues.
\item $\tr(A)=\sum\lambda_i$ (sum of eigenvalues, with multiplicity).
\item $\det(A)=\prod\lambda_i$ (product of eigenvalues, with multiplicity).
\end{itemize}

\subsect{Algebraic vs.\ Geometric Multiplicity}
\begin{itemize}[nosep,leftmargin=*]
\item Algebraic mult.\ of $\lambda=$ multiplicity as root of $p(\lambda)$.
\item Geometric mult.\ of $\lambda=\dim E_\lambda=\nullity(A-\lambda I)$.
\item Always: $1\le\text{geo}\le\text{alg}$.
\end{itemize}

\subsect{Key Facts}
\begin{itemize}[nosep,leftmargin=*]
\item Eigenvectors for distinct eigenvalues are linearly independent.
\item Similar matrices ($B=P^{-1}AP$) have the same eigenvalues, trace, determinant, and characteristic polynomial.
\item If $\lambda$ is an eigenvalue of $A$: eigenvalue of $A^k$ is $\lambda^k$; of $A^{-1}$ is $1/\lambda$; of $A+cI$ is $\lambda+c$.
\end{itemize}

%======================================================================
\sect{6. Diagonalization}
%======================================================================

\subsect{When Is $A$ Diagonalizable?}
$A$ is diagonalizable $\iff$ $\exists$ basis of eigenvectors $\iff$ geo mult $=$ alg mult for every eigenvalue.
\begin{itemize}[nosep,leftmargin=*]
\item $n$ distinct eigenvalues $\Rightarrow$ diagonalizable (sufficient, not necessary).
\item $A=PDP^{-1}$ where $D=\diag(\lambda_1,\ldots,\lambda_n)$, columns of $P$ are eigenvectors.
\item Real symmetric matrices are always diagonalizable (and orthogonally diagonalizable).
\end{itemize}

\subsect{Cayley--Hamilton Theorem}
Every matrix satisfies its own characteristic polynomial: $p(A)=0$.\\
\textit{Example:} If $p(\lambda)=\lambda^2-5\lambda+6$, then $A^2-5A+6I=0$.

%======================================================================
\sect{7. Inner Product Spaces}
%======================================================================

\subsect{Inner Product on $\mathbb{R}^n$}
$\langle u,v\rangle = u\cdot v = u^Tv = \sum u_iv_i$.\\
$\|v\|=\sqrt{\langle v,v\rangle}$.

\subsect{Orthogonality}
\begin{itemize}[nosep,leftmargin=*]
\item $u\perp v \iff \langle u,v\rangle=0$.
\item Orthogonal set of nonzero vectors is linearly independent.
\item \textbf{Gram--Schmidt:} produces orthogonal (or orthonormal) basis from any basis.
\item Orthogonal complement: $W^\perp=\{v:\langle v,w\rangle=0\;\forall w\in W\}$.
\item $\dim W+\dim W^\perp=\dim V$.
\item $V=W\oplus W^\perp$.
\end{itemize}

\subsect{Orthogonal \& Symmetric Matrices}
\begin{itemize}[nosep,leftmargin=*]
\item \textbf{Orthogonal:} $Q^TQ=I$ $\iff$ columns form an orthonormal basis.\\
  $\det Q=\pm 1$, preserves lengths and angles.
\item \textbf{Symmetric:} $A^T=A$.\\
  All eigenvalues are real.\\
  Eigenvectors for distinct eigenvalues are orthogonal.\\
  $A=QDQ^T$ (spectral theorem) with $Q$ orthogonal.
\end{itemize}

\subsect{Positive Definite Matrices}
Symmetric $A$ is positive definite if $x^TAx>0$ for all $x\neq 0$.\\
Equivalent: all eigenvalues $>0$; all leading principal minors $>0$.

%======================================================================
\sect{8. Special Matrices \& Decompositions}
%======================================================================

\subsect{Special Types}
\begin{itemize}[nosep,leftmargin=*]
\item \textbf{Nilpotent:} $A^k=0$ for some $k$. All eigenvalues are 0.\ $\det A=0$, $\tr A=0$.
\item \textbf{Idempotent:} $A^2=A$. Eigenvalues $\in\{0,1\}$.\ $\rank(A)=\tr(A)$.
\item \textbf{Projection:} idempotent. Orthogonal projection if also symmetric.
\item \textbf{Skew-symmetric:} $A^T=-A$. Diagonal entries are 0. Eigenvalues are purely imaginary (or 0). $\det A\ge 0$ if $n$ is even.
\end{itemize}

\subsect{Jordan Normal Form (Summary)}
Every matrix over $\mathbb{C}$ is similar to a Jordan matrix: block diagonal with Jordan blocks $J_k(\lambda)=\lambda I_k + N_k$ (upper shift). Number of Jordan blocks for $\lambda=$ geometric multiplicity.

%======================================================================
\sect{9. Quick Formulas \& Tricks}
%======================================================================

\subsect{Trace \& Determinant}
\begin{itemize}[nosep,leftmargin=*]
\item $\tr(A+B)=\tr A+\tr B$, \quad $\tr(cA)=c\cdot\tr A$.
\item $\tr(AB)=\tr(BA)$ (even if $AB\neq BA$).
\item $\tr(A)=\sum\lambda_i$, \quad $\det(A)=\prod\lambda_i$.
\end{itemize}

\subsect{$2\times 2$ Eigenvalue Shortcut}
For $A=\begin{pmatrix}a&b\\c&d\end{pmatrix}$:
\[
\lambda = \frac{\tr A\pm\sqrt{(\tr A)^2-4\det A}}{2}.
\]

\subsect{Rank Inequalities}
\begin{itemize}[nosep,leftmargin=*]
\item $\rank(A+B)\le\rank A+\rank B$.
\item $\rank(AB)\ge\rank A+\rank B-n$ (Sylvester).
\item $\rank(A)=\rank(A^T)=\rank(A^TA)$.
\end{itemize}

\subsect{Change of Basis}
If $P$ is the change-of-basis matrix from $B'$ to $B$, then:
$[T]_{B'}=P^{-1}[T]_B\,P$.

\bigskip
\hrule
\smallskip
\begin{center}
\textit{Key chain:}\; $\det A\neq 0$ $\iff$ invertible $\iff$ $\rank=n$ $\iff$ $\ker=\{0\}$ $\iff$ all eigenvalues $\neq 0$.
\end{center}

\end{document}
